\chapter{基于梯度子空间融合的无回放联邦持续医学图像分类研究}
\label{chap:fedsubmerge}

\section{引言}
\label{sec:fedsubmerge-introduction}

\section{问题定义与研究动机}
\label{sec:fedsubmerge-motivation}

\section{无回放联邦持续学习问题建模}
\label{sec:fedsubmerge-formulation}

\section{FedSubMerge 方法}
\label{sec:fedsubmerge-method}

\subsection{方法总体框架}
\label{subsec:fedsubmerge-overview}

\subsection{客户端主梯度子空间构建}
\label{subsec:fedsubmerge-client-pgs}

\subsection{服务器端跨客户端子空间融合}
\label{subsec:fedsubmerge-server-merge}

\subsection{基于正交补投影的本地训练}
\label{subsec:fedsubmerge-local-projection}

\subsection{自适应逐层客户端选择}
\label{subsec:fedsubmerge-adaptive-selection}

\subsection{客户端特定保护子空间构建}
\label{subsec:fedsubmerge-personalized-subspace}

\subsection{算法流程与实现细节}
\label{subsec:fedsubmerge-algorithm}

\section{方法理论分析}
\label{sec:fedsubmerge-theory}

\section{实验设计与结果分析}
\label{sec:fedsubmerge-experiments}

\section{方法讨论}
\label{sec:fedsubmerge-discussion}
% 讨论历史访问、异质性、通信对象、计算与存储开销、任务边界及隐私表述边界。

\section{本章小结}
\label{sec:fedsubmerge-summary}
